% §894 makes **one** hyphenation attempt per glue node, and §866 is the only
% thing that calls it:
%
%     When the following code is activated, the |line_break| procedure is
%     in its second pass, and |cur_p| points to a glue node.
%
% §896 then walks forward looking for the word to try. A character whose
% `\lccode` is zero — an opening paren, a quote — is *stepped over* on the way:
%
%     if lc_code(c)<>0 then
%       if (lc_code(c)=c)or(uc_hyph>0) then goto done2
%       else goto done1;
%     continue: prev_s:=s; s:=link(prev_s);
%
% so `(knowledge)` hyphenates as `(knowl-edge)`. §897 then collects letters and
% stops at the first character with a zero `\lccode`:
%
%     if lc_code(c)=0 then goto done3;
%
% and that is the end of the attempt. There is no path back to §896 for whatever
% follows, so a second word behind the same glue is never tried.
%
% `(semi)stable` is the case that separates the two halves. The word §896 finds
% is `semi`; four letters cannot be hyphenated with `\lefthyphenmin=2` and
% `\righthyphenmin=3`; and `stable` — which on its own is `sta-ble` — is never
% looked at. Stepping over the `)` as though still searching gives
% `(semi)sta-ble`, one discretionary that pdfTeX does not have, and it was worth
% a whole word of page 2 in 2608.17397v1.
\input plain
\hsize=80pt
\parindent=0pt
\pretolerance=-1
\tolerance=9999
\showboxbreadth=100
\showboxdepth=10
% The word behind the second `)` must stay whole.
\setbox0=\vbox{\noindent XX (semi)stable YY\par}
\showbox0
% …while a word behind a *leading* zero-lccode character is still reached.
\setbox0=\vbox{\noindent XX (knowledge) YY\par}
\showbox0
% Two of them, and a quote, ahead of the word.
\setbox0=\vbox{\noindent XX ((knowledge YY\par}
\showbox0
\setbox0=\vbox{\noindent XX ``knowledge YY\par}
\showbox0
% No punctuation at all: the plain case still hyphenates.
\setbox0=\vbox{\noindent XX knowledge YY\par}
\showbox0
% A word, then punctuation, then a second word — the second gets no attempt.
\setbox0=\vbox{\noindent XX able)knowledge YY\par}
\showbox0
% …but a glue between them arms a fresh one.
\setbox0=\vbox{\noindent XX able) knowledge YY\par}
\showbox0
% The first word of a paragraph is preceded by the `\parindent` box, not glue,
% so it is never hyphenated — `\noindent` above makes that box empty but it is
% still not a glue node.
\setbox0=\vbox{\noindent knowledge XX\par}
\showbox0
\end
